\section{Semi-leptonic decay channels (Day 1, morning)}

\subsection{Experimental review for diboson and VBS}
\speaker{Ankita Mehta}{Ghent University, for the ATLAS and CMS Collaborations}
The session opened with a comprehensive review of the LHC electroweak
multiboson program at the precision and energy frontiers. Diboson production
($WW$, $WZ$, $ZZ$, $V\gamma$) has been measured inclusively and differentially
at 5.02, 13, and 13.6~TeV, with precision that in many areas already exceeds
the original projections, thanks to detector performance, advanced analysis
techniques, and higher-order theory (up to NNLO~QCD + NLO~EW). Highlights
included the CMS anomalous-coupling analysis of semi-leptonic $WV$ with
ParticleNet-tagged boosted bosons (currently among the strongest bounds on
triple-gauge and four-fermion operators), the ATLAS $W\gamma$ measurement with
radiation-amplitude-zero and polarization-sensitive observables, and the ATLAS
$ZZ\to 2\ell 2\nu(jj)$ cross sections with aTGC/aQGC constraints. On the VBS
side, all major channels ($W^{\pm}W^{\pm}$, $WZ$, $ZZ$) are now observed; the
first Run-3 VBS measurement (CMS $W^{\pm}W^{\pm}$ and $WZ$ at 13.6~TeV,
including first differential $WZ$ distributions) is statistically limited, and
semi-leptonic channels ($WVjj$, $ZVjj$, $VVjj$) increasingly rely on
DNN/RNN-based signal--background discrimination. The speaker stressed that
semi-leptonic modes are challenging but that ML-based boson tagging has been a
game changer, that VBS simulations are still LO in many cases, and that
combinations and SMEFT interpretations (including the ATLAS EW global fit) make
these measurements key inputs to global searches for new physics.

\subsection{Theory review on polarised bosons}
\speaker{Rene Poncelet}{IFJ PAN Krak\'ow}
The theory review traced the path from the motivation (unitarity of
$W_LW_L$ scattering and its link to EWSB) to the state of the art in polarized
predictions. Since boson polarization is not directly observable, decay
products act as polarimeters; the modern approach defines polarized cross
sections for on-shell bosons in the double-pole or narrow-width approximation
(DPA/NWA), generating separately polarized samples that are fitted to data,
with interference contributions kept under control. Higher-order corrections
are essential: polarization fractions are shape-dependent and the corrections
themselves depend on the polarization state, so unpolarized $K$-factors
distort spectra; NNLO~QCD is needed for percent-level scale uncertainties
(illustrated with polarized $W$+jet at NNLO and mock-data fits). The talk
reviewed the rapidly growing literature of polarized calculations---NLO~QCD/EW
for $WW$, $WZ$, $ZZ$, NNLO~QCD for $WW$ and (new) NNLO~QCD+NLO~EW for $WZ$ with
theory-nuisance-parameter uncertainties, semi-leptonic decays at NLO~QCD
(resolved and boosted setups, where $LL$ turns out to be rather insensitive to
QCD corrections while $TT$/mixed states receive large ones)---as well as
Monte Carlo developments (polarized MadGraph, SHERPA nLO+PS, POWHEG-BOX-RES
NLO+PS, and SMEFT@NLO+PS). A dedicated COMETA study of polarized $ZZ$
production was presented, comparing NWA/DPA/off-shell treatments,
matching/merging schemes, and best predictions against the ATLAS polarization
fractions (agreement, dominated by experimental uncertainties). The outlook
identified hadronic polarization taggers, higher-order corrections to decays,
and parton-shower matching as the next frontiers.

\subsection{Gauge boson polarization as a probe of SMEFT effects in $WZ$ production (remote)}
\speaker{Maryam Bayat Makou}{Southern Methodist University; with S.~Bhattacharya}
This contribution presented a feasibility study, carried out within COMETA, of
the interplay between dimension-6 SMEFT operators and gauge boson polarization
in $WZ\to 3\ell$ production at $\sqrt{s}=13$~TeV. Dedicated polarized samples
($W_TZ_T$, $W_LZ_L$, $W_TZ_L$, $W_LZ_T$; 100k events each, MadGraph + Pythia +
Delphes) were compared between SM and SMEFT using a broad set of kinematic and
angular observables. EFT effects grow in the high-$\pT$ tails (clearly visible
in the leading-lepton $\pT$ above $\sim$200~GeV, dominated by the
energy-growing behavior of $C_W$); ranking the operators by a
polarization-shape metric shows that $c_{HWB}$, $c_{HDD}$ and $c_{Hl3}$ produce
the largest genuine distortions of the polarization composition, while others
(e.g.\ $c_{HWB}$ at large coefficient) act mainly on the overall rate. A DNN
trained on SM polarized samples to separate longitudinal from transverse
configurations learns physically meaningful observables (the top feature being
the lepton rapidity separation $\Delta y$), and EFT events systematically shift
the learned polarization scores, providing operator-dependent ``polarization
fingerprints''. A two-output architecture predicting $P(W_L)$ and $P(Z_L)$
separately was advocated as more physical and better suited for EFT studies. A
publication is in preparation.

\subsection{Discussion}
The session closed with a discussion on the complementarity between
polarization observables and traditional EFT analyses, on the treatment of
interference and off-shell effects in polarized templates, and on the
priorities for extending these studies to semi-leptonic final states.
